\chapter{Вывод}
\label{ch:chap6}

В ходе выполнения лабораторной работы была составлена модель цифровой САУ температуры. Настроены параметры дискретного регулятора для нескольких случаев: при разных периодах дискретизации и отклонениях в большую и меньшую сторону постоянной времени $T_2$ объекта управления. 

Исследовано поведение замкнутой системы при ступенчатом изменении задающего воздействия, при ступенчатом изменении возмущающего воздействия и при возмущающем воздействии, изменяющемся по случайному закону. В случае изменения задающего воздействия скачком система через определенное время вновь становится устойчивой относительно нового значения задающего воздействия. То же справедливо и для скачкообразного изменения возмущающего воздействия. В случае случайного возмущающего воздействия на систему, система стабилизируется около целевого значения $g(t)$ и остается ограниченным в условиях ограниченного возмущения.

При уменьшении времени дискретизации уменьшаются амплитуды переходных процессов, система быстрее реагирует на изменения возмущающего воздействия. Переходные процессы становятся апериодическими.

При уменьшении $T_2$ на $20 \%$ увеличивается амплитуда переходных процессов, также повышается частота колебаний в процессе стабилизации системы.

\endinput